% paper_16_periodicity.tex — GeoVac Paper 16
% Chemical Periodicity as S_N Representation Theory on S^{3N-1}
% Status: Draft, March 18, 2026
% Target: Journal of Chemical Physics or Physical Review Letters
\documentclass[aps,prl,twocolumn,superscriptaddress,longbibliography]{revtex4-2}

\usepackage{amsmath,amssymb,graphicx,xcolor,bm,braket,dcolumn,booktabs}
\usepackage{hyperref}

\newcommand{\Nelectron}{N}
\newcommand{\mufree}{\mu_{\mathrm{free}}}
\newcommand{\neff}{n_{\mathrm{eff}}}
\newcommand{\leff}{l_{\mathrm{eff}}}
\newcommand{\Zeff}{Z_{\mathrm{eff}}}
\newcommand{\Ha}{\,\mathrm{Ha}}
\newcommand{\eV}{\,\mathrm{eV}}

\begin{document}

\title{The Periodic Table Recast as $S_N$ Representation Theory\\
on $S^{3N-1}$}
\thanks{Paper 16 in the Geometric Vacuum series}

\author{J.~Loutey}
\affiliation{Independent Researcher, Kent, Washington}

\date{March 18, 2026}

%% ========================================================================
%% ABSTRACT
%% ========================================================================

\begin{abstract}
The non-relativistic $N$-electron Schr\"odinger equation, after
separation of the hyperradius, reduces to angular dynamics on the
hypersphere $S^{3N-1}$.  The angular kinetic energy is the
$\mathrm{SO}(3N)$ Casimir operator with eigenvalue
$\mufree = \nu(\nu + 3N - 2)/2$.  Fermionic antisymmetry restricts
the spatial wavefunction to a specific $S_N$ irreducible
representation---the conjugate of the spin Young diagram---whose
first-column length $\lambda_1$ fixes the minimum angular quantum
number $\nu = N - \lambda_1$.  For all atoms with ground-state spin
$S < N/2$, the spatial irrep has $\lambda_1 = 2$, giving the
universal result $\nu = N - 2$ and $\mufree = 2(N-2)^2$.  This
formula, combined with the structure of the $S_N$ irrep, yields a
classification of atoms into four types---A (trivial), B~(noble gas),
C~(alkali), D~(alkaline earth)---that recasts the periodic table's
group structure in representation-theoretic language.  The periodic law
corresponds to the repetition of the $S_N$ irrep sequence
$\mathrm{C} \to \mathrm{D} \to \cdots \to \mathrm{B}$ within each
period.  This classification maps onto, rather than predicts, the
known periodic table structure; its value lies in providing a
unified group-theoretic framework for a pattern that is usually
described empirically.  In the per-pair limit, $\mufree$
contributes exactly~4 units of angular momentum per antisymmetric
electron pair as $N \to \infty$.  We show that the topological
structure ($\nu$, $\mufree$) is smooth through $Z = 1/\alpha \approx
137$ and that the Dirac instability is a \emph{metric} singularity on
$S^{3N-1}$, not a topological one.
\end{abstract}

\maketitle

%% ========================================================================
%% I. INTRODUCTION
%% ========================================================================

\section{Introduction}
\label{sec:intro}

The periodic table of elements, established empirically by Mendeleev
in 1869~\cite{Mendeleev1869}, received its first theoretical
justification from quantum mechanics through the shell model and the
aufbau principle.  The $2n^2$ orbital counting rule, the grouping of
elements by valence shell, and the progression of ionization energies
across a period are all consequences of the hydrogen-like
Schr\"odinger equation and its multi-electron extensions.

Despite this success, the periodic law itself---the recurrence of
chemical properties with increasing atomic number---has not been
derived from a single mathematical principle.  Shell-filling rules are
empirical (the Madelung rule has exceptions), Hund's rules require
separate justification, and the distinction between noble gases, alkali
metals, and transition metals is described rather than predicted.

In this paper, we show that the periodic table's group structure
can be recast in the language of $S_N$ representation theory acting
on the spatial part of the $N$-electron wavefunction on the
hypersphere $S^{3N-1}$.  The key results are:

\begin{enumerate}
  \item The ground-state angular quantum number is universally $\nu =
    N - 2$ for atoms with spin $S < N/2$ (Sec.~\ref{sec:framework}).
  \item The Pauli centrifugal cost $\mufree = 2(N-2)^2$ grows
    quadratically, contributing exactly~4 units per electron pair in
    the large-$N$ limit (Sec.~\ref{sec:mu_formula}).
  \item The shape of the $S_N$ irrep classifies atoms into chemical
    types that recast the group structure of the periodic table in
    representation-theoretic language
    (Sec.~\ref{sec:classification}).
  \item The periodic law corresponds to the repetition of $S_N$ irrep
    sequences (Sec.~\ref{sec:periodic_law}).
  \item The Dirac instability at $Z\alpha \to 1$ is a metric
    singularity, not a topological one; the representation theory
    is smooth through $Z = 137$
    (Sec.~\ref{sec:topology_vs_metric}).
\end{enumerate}

%% ========================================================================
%% I-B. WHAT IS NEW VS. WHAT IS KNOWN
%% ========================================================================

\section{What is New vs.\ What is Known}
\label{sec:new_vs_known}

Several ingredients of the present analysis are standard; it is
important to distinguish them from the new contributions.

\subsection{Known results}

\begin{itemize}
  \item \textbf{Hyperspherical formulation.}  The separation of the
    $N$-electron Schr\"odinger equation into hyperradial and angular
    parts on $S^{3N-1}$, and the identification of the angular kinetic
    energy with the $\mathrm{SO}(3N)$ Casimir, are
    established~\cite{Avery2000, Fock1954}.

  \item \textbf{$S_N$ representation theory for fermions.}  The use of
    Young diagrams and Schur--Weyl duality to classify spatial--spin
    factorizations of antisymmetric wavefunctions is textbook
    material~\cite{Weyl1946, FultonHarris1991}.

  \item \textbf{The Casimir eigenvalue formula.}  The expression
    $\mufree = \nu(\nu + 3N - 2)/2$ follows directly from the
    standard $\mathrm{SO}(3N)$ representation theory.

  \item \textbf{The periodic table itself.}  The classification of
    elements into noble gases, alkali metals, alkaline earths, and
    transition metals is an empirical fact that guided the
    construction presented here.
\end{itemize}

\subsection{New contributions}

\begin{enumerate}
  \item \textbf{The universal $\nu = N - 2$ result.}  Theorem~1 shows
    that, for all ground-state atoms with $S < N/2$, the spatial
    irrep has first-column length $\lambda_1 = 2$, yielding the
    universal angular quantum number $\nu = N - 2$ and $\mufree =
    2(N-2)^2$.  This specific identification appears to be new.

  \item \textbf{The four-type classification (A/B/C/D).}  The
    observation that the $S_N$ irrep shape naturally partitions atoms
    into types that map onto the periodic table's group structure
    provides a group-theoretic reformulation of the empirical pattern.

  \item \textbf{The per-pair angular momentum limit.}  The result that
    $\mufree / \binom{N}{2} \to 4$ as $N \to \infty$---each
    antisymmetric electron pair contributes exactly~4 units of
    angular momentum---connects the periodic pattern to a
    quantitative angular momentum budget.

  \item \textbf{The Dirac instability as metric vs.\ topological
    singularity.}  The analysis showing that $\nu$, $\mufree$, and
    the $S_N$ irrep are smooth through $Z = 137$ while the metric
    diverges clarifies the nature of the upper limit on the periodic
    table.
\end{enumerate}

%% ========================================================================
%% III. THE HYPERSPHERICAL FRAMEWORK
%% ========================================================================

\section{The Hyperspherical Framework}
\label{sec:framework}

\subsection{Configuration space and angular dynamics}
\label{sec:config}

Consider $N$ electrons in a central potential $V(r_1, \ldots, r_N)$.
After extracting the center-of-mass motion, the configuration space is
$\mathbb{R}^{3N}$.  Introducing hyperspherical coordinates
$(R, \hat{\Omega})$ where $R = \sqrt{\sum_i r_i^2}$ is the
hyperradius and $\hat{\Omega} \in S^{3N-1}$ collects $3N-1$ angular
variables, the kinetic energy separates as~\cite{Avery2000,
Fock1954}:
\begin{equation}
  T = -\frac{1}{2}\left(\frac{\partial^2}{\partial R^2}
    + \frac{3N-1}{R}\frac{\partial}{\partial R}
    - \frac{\Lambda^2}{R^2}\right),
  \label{eq:kinetic}
\end{equation}
where $\Lambda^2$ is the grand angular momentum operator---the
Laplace--Beltrami operator on $S^{3N-1}$.  Its eigenvalues are the
$\mathrm{SO}(3N)$ Casimir invariants:
\begin{equation}
  \Lambda^2 Y_\nu = \nu(\nu + 3N - 2)\, Y_\nu \equiv 2\mufree\, Y_\nu,
  \label{eq:casimir}
\end{equation}
with $\nu = 0, 1, 2, \ldots$ and hyperspherical harmonics $Y_\nu$.
The effective angular momentum quantum numbers are
\begin{equation}
  \leff = \frac{3N - 3}{2}, \quad \neff = \frac{3N - 1}{2}.
  \label{eq:eff_qn}
\end{equation}

The centrifugal barrier in the hyperradial equation is
$\leff(\leff + 1)/R^2 = (3N-3)(3N-1)/(4R^2)$, showing that even in
the lowest angular channel, $N$-electron atoms face a centrifugal cost
that grows as $\sim 9N^2/4$.

\subsection{Fermionic antisymmetry and $S_N$ irreps}
\label{sec:sn_irreps}

The total wavefunction must be antisymmetric under $S_N$ (permutation
of electrons).  Since the total state factorizes as $|\Psi\rangle =
|\text{spatial}\rangle \otimes |\text{spin}\rangle$, the spatial and
spin parts must carry conjugate $S_N$ irreducible representations.

For $N$ spin-1/2 fermions with total spin $S$, the spin irrep is a
Young diagram with at most two rows: $[n_\uparrow, n_\downarrow]$
where $n_\uparrow = N/2 + S$ and $n_\downarrow = N/2 - S$.  By
Schur--Weyl duality~\cite{Weyl1946, FultonHarris1991}, the spatial
irrep is the \emph{conjugate} (transposed) Young diagram.

\noindent\textbf{Theorem 1} (Universal $\lambda_1$).
\label{thm:lambda1}
\textit{For any $N$-electron atom with ground-state spin $S < N/2$, the
spatial Young diagram has first-column length $\lambda_1 = 2$.}

\medskip\noindent\textit{Proof.}
The spin diagram has two rows: $[n_\uparrow, n_\downarrow]$ with
$n_\downarrow \geq 1$ (since $S < N/2$ implies $n_\downarrow > 0$).
The conjugate diagram has $n_\downarrow$ columns of height~2 and
$n_\uparrow - n_\downarrow = 2S$ columns of height~1.  Therefore the
first column has height $\lambda_1 = 2$, independent of $S$ (provided
$S < N/2$). \hfill$\square$

\subsection{The ground-state angular quantum number}
\label{sec:nu_gs}

The minimum $\nu$ compatible with a given $S_N$ irrep is determined by
the constraint that the hyperspherical harmonic $Y_\nu$ must transform
in that irrep under particle permutations.  For the spatial irrep with
first-column length $\lambda_1$, the minimum angular excitation
is~\cite{Avery2000}:
\begin{equation}
  \nu_{\min} = N - \lambda_1.
  \label{eq:nu_min}
\end{equation}

By Theorem~\ref{thm:lambda1}, all atoms with $S < N/2$ have
\begin{equation}
  \boxed{\nu = N - 2}
  \label{eq:nu_universal}
\end{equation}
as their ground-state angular quantum number.  This is a
\emph{universal} result: it depends only on the electron count $N$,
not on the nuclear charge $Z$, the spin $S$ (provided $S < N/2$), or
any details of the electron configuration.

The sole exception is the maximally polarized state $S = N/2$ (all
spins parallel), which has $\lambda_1 = N$ and $\nu = 0$.  This state
is completely symmetric in space and carries zero angular momentum---it
is the Laughlin-like state with no Pauli nodes.

%% ========================================================================
%% III. THE mu_free FORMULA
%% ========================================================================

\section{The $\mufree$ Formula}
\label{sec:mu_formula}

Substituting $\nu = N - 2$ into Eq.~(\ref{eq:casimir}):
\begin{align}
  \mufree &= \frac{\nu(\nu + 3N - 2)}{2}
           = \frac{(N-2)(N - 2 + 3N - 2)}{2} \notag \\
          &= \frac{(N-2)(4N-4)}{2}
           = \boxed{2(N-2)^2}.
  \label{eq:mu_formula}
\end{align}

This is the \emph{Pauli centrifugal cost}: the minimum angular kinetic
energy imposed by fermionic antisymmetry on $S^{3N-1}$.

\subsection{Asymptotic behavior}
\label{sec:asymptotic}

Several limits of Eq.~(\ref{eq:mu_formula}) are physically
informative:

\begin{enumerate}
  \item \textbf{Per-electron scaling.}  $\mufree / N^2 = 2 - 8/N +
    8/N^2 \to 2$ as $N \to \infty$.

  \item \textbf{Per-pair scaling.}  The number of electron pairs is
    $\binom{N}{2} = N(N-1)/2$.  The Pauli cost per pair is
    \begin{equation}
      \frac{\mufree}{\binom{N}{2}}
        = \frac{4(N-2)^2}{N(N-1)} \longrightarrow 4
        \quad (N \to \infty).
      \label{eq:per_pair}
    \end{equation}
    Each antisymmetric electron pair contributes exactly~4 units
    of angular momentum to the ground-state Casimir.

  \item \textbf{Excitation fraction.}  The ratio $\nu/(3N-2) =
    (N-2)/(3N-2) \to 1/3$ monotonically from below.  The ground state
    always sits at approximately one-third of the natural angular
    scale of the sphere, regardless of $N$.

  \item \textbf{Fractional ionization cost.}  Removing one electron
    changes the sphere dimension by $3/(3N-1) \to 0$.  For heavy
    atoms, ionization is a \emph{small perturbation} to the angular
    topology.
\end{enumerate}

Table~\ref{tab:mu_values} collects these quantities for atoms across
the periodic table.

\begin{table}[h]
\centering
\caption{Ground-state hyperspherical parameters for neutral atoms
($S < N/2$).  $\nu = N-2$, $\mufree = 2(N-2)^2$, and the asymptotic
ratios $\mufree/N^2 \to 2$, $\mufree/\binom{N}{2} \to 4$,
$\nu/(3N-2) \to 1/3$.}
\label{tab:mu_values}
\begin{tabular}{lrrrrrc}
\toprule
Atom & $N$ & $S^{3N-1}$ & $\nu$ & $\mufree$ & $\mufree/N^2$ & Type \\
\midrule
H     &   1 & $S^2$    &   0 &       0 & 0.00 & A \\
He    &   2 & $S^5$    &   0 &       0 & 0.00 & B \\
Li    &   3 & $S^8$    &   1 &       2 & 0.22 & C \\
Be    &   4 & $S^{11}$ &   2 &       8 & 0.50 & D \\
N     &   7 & $S^{20}$ &   5 &      50 & 1.02 & E \\
Ne    &  10 & $S^{29}$ &   8 &     128 & 1.28 & B \\
Na    &  11 & $S^{32}$ &   9 &     162 & 1.34 & C \\
Mg    &  12 & $S^{35}$ &  10 &     200 & 1.39 & D \\
Ar    &  18 & $S^{53}$ &  16 &     512 & 1.58 & B \\
Kr    &  36 & $S^{107}$&  34 &   2312 & 1.78 & B \\
Xe    &  54 & $S^{161}$&  52 &   5408 & 1.85 & B \\
Rn    &  86 & $S^{257}$&  84 &  14112 & 1.91 & B \\
Og    & 118 & $S^{353}$& 116 &  26912 & 1.93 & B \\
\bottomrule
\end{tabular}
\end{table}

%% ========================================================================
%% IV. STRUCTURE TYPE CLASSIFICATION
%% ========================================================================

\section{Structure Type Classification}
\label{sec:classification}

While $\nu = N - 2$ and $\mufree = 2(N-2)^2$ are universal, the
\emph{shape} of the $S_N$ irrep varies with the electron
configuration.  This shape encodes the internal structure of the
atom---how the $N$~electrons partition into core and valence.

We define four structure types based on the spatial Young diagram:

\subsection{Type A: Trivial ($N = 1$)}

The hydrogen atom has a single electron.  The trivial irrep $[1]$ of
$S_1$ gives $\lambda_1 = 1$, $\nu = 0$, $\mufree = 0$.  There is no
Pauli constraint.  The wavefunction lives on $S^2$ (a single
electron's angular space).

\subsection{Type B: Noble gas (democratic structure)}

Noble gases have closed-shell ground states with $S = 0$.  The spin
diagram is $[N/2,\, N/2]$ and the spatial irrep is its conjugate:
$[2, 2, \ldots, 2]$ ($N/2$ rows of length~2).

This is the most \emph{democratic} irrep: all columns have the same
height.  No electron is distinguished.  The wavefunction has the full
point-group symmetry of the Young diagram, and ionization requires
breaking this symmetry---a costly operation reflected in the high
ionization energies of noble gases (He:~$24.6\eV$, Ne:~$21.6\eV$,
Ar:~$15.8\eV$).

\subsection{Type C: Alkali metal (hierarchical structure)}

Alkali metals have one valence electron outside a noble-gas core, with
$S = 1/2$.  The spin diagram is $[(N+1)/2, \,(N-1)/2]$ and the
spatial irrep is $[2, 2, \ldots, 2, 1]$ ($= [2^{(N-1)/2}, 1]$).

This is a \emph{hierarchical} irrep: one column (the valence electron)
is shorter than the rest.  The wavefunction has a natural
decomposition into a symmetric core and a weakly bound valence
electron.  Ionization removes the short column---a topologically clean
surgery on $S^{3N-1}$:
\begin{equation}
  S^{3N-1} \to S^{3(N-1)-1} \times \mathbb{R}.
  \label{eq:ionization_surgery}
\end{equation}
The low cost of this surgery explains the low ionization energies
of alkali metals (Li:~$5.4\eV$, Na:~$5.1\eV$, K:~$4.3\eV$).

\subsection{Type D: Alkaline earth (core + pair)}

Alkaline earth metals have two valence electrons outside a noble-gas
core, with $S = 0$.  The spatial irrep is $[2, 2, \ldots, 2]$ (same
shape as Type~B) but the \emph{occupancy pattern} distinguishes the
core from the two valence electrons.

The ionization energy is intermediate (Be:~$9.3\eV$, Mg:~$7.6\eV$,
Ca:~$6.1\eV$)---higher than alkali metals (the valence pair is harder
to break) but lower than noble gases (the valence pair is still
distinguishable from the core).

Table~\ref{tab:types} summarizes the classification.

\begin{table}[h]
\centering
\caption{Structure type classification.  The spatial Young diagram is
the conjugate of the spin diagram.  IE$_1$ is the first ionization
energy (experimental).}
\label{tab:types}
\begin{tabular}{cllcr}
\toprule
Type & Spatial irrep & Chemistry & Examples & IE$_1$ range \\
\midrule
A & $[1]$       & Monatomic       & H            & $13.6\eV$ \\
B & $[2^{N/2}]$ & Noble gas       & He, Ne, Ar   & $15$--$25\eV$ \\
C & $[2^{(N-1)/2}, 1]$ & Alkali  & Li, Na, K    & $4$--$5\eV$ \\
D & $[2^{N/2}]$ & Alkaline earth  & Be, Mg, Ca   & $6$--$9\eV$ \\
E & Various     & Open shell      & B--F, d-block& $8$--$17\eV$ \\
\bottomrule
\end{tabular}
\end{table}

%% ========================================================================
%% V. THE PERIODIC LAW AS IRREP SEQUENCE
%% ========================================================================

\section{The Periodic Law as Irrep Sequence}
\label{sec:periodic_law}

The periodic table's group structure arises because the sequence of
$S_N$ structure types \emph{repeats} with each period.  Within a given
period (defined by the outermost principal quantum number $n$), the
element sequence is:
\begin{equation}
  \mathrm{C} \to \mathrm{D} \to \underbrace{\mathrm{E} \to \cdots \to
  \mathrm{E}}_{p\text{- and }d\text{-block}} \to \mathrm{B}.
  \label{eq:period_sequence}
\end{equation}
Each period begins with an alkali metal (Type~C: one electron in the
new shell) and ends with a noble gas (Type~B: closed shell).  The
intermediate elements (Type~E) have partially filled subshells with
varying spin and Young diagram shapes.

In group-theoretic language, the periodic law can be understood as
follows:
\begin{quote}
\emph{The chemical properties of elements are periodic functions of
atomic number because the sequence of $S_N$ spatial irreps repeats
whenever a new principal shell begins filling.}
\end{quote}
This is a reformulation, not a derivation from first principles: the
shell-filling order itself (which determines where the sequence
repeats) is an empirical input, not a consequence of $S_N$
representation theory alone.

The key distinction between Types~B and~D is worth emphasizing.
Both have the same Young diagram shape $[2^{N/2}]$, but they differ
in the \emph{occupancy pattern}: Type~B fills all available subshells,
while Type~D has two electrons in a new shell above a closed core.
This distinction is invisible at Level~1 (representation theory) and
becomes visible only at Level~2 (the perturbative energy, which
depends on $Z$, orbital angular momenta, and electron--electron
repulsion).

\subsection{Isoelectronic series}
\label{sec:isoelectronic}

Atoms with the same $N$ but different $Z$ share identical $S_N$ irreps
and structure types.  The 10-electron series
\begin{equation}
  \text{Ne},\ \text{Na}^+,\ \text{Mg}^{2+},\ \text{Al}^{3+},\ \ldots
\end{equation}
all have the Type~B spatial irrep $[2,2,2,2,2]$ with $\nu = 8$ and
$\mufree = 128$.  Their chemistry differs only through the nuclear
charge~$Z$, which enters at Level~2.  This cleanly separates
\emph{topology} (same for the series) from \emph{potential} (varies
with $Z$).

\subsection{Topology flow: ionization as surgery}
\label{sec:topology_flow}

Ionization changes the electron count $N \to N - 1$, transforming the
configuration space via Eq.~(\ref{eq:ionization_surgery}).  This is a
codimension-3 surgery that removes three dimensions (one electron's
$\mathbb{R}^3$) from the hypersphere.  The fractional dimension loss
is $3/(3N-1)$, which decreases monotonically:
\begin{equation}
  \frac{3}{3N - 1} \to 0 \quad (N \to \infty).
\end{equation}
For heavy atoms, ionization is a negligible topological perturbation.
This correlates with the empirical observation that ionization energies
decrease down a group: removing one electron from $S^{353}$ (Og) is
geometrically less disruptive than removing one from $S^5$ (He).

%% ========================================================================
%% VI. TWO-LEVEL SEPARATION
%% ========================================================================

\section{Topology versus Metric: Two-Level Structure}
\label{sec:topology_vs_metric}

The preceding analysis reveals a clean two-level separation:

\begin{itemize}
  \item \textbf{Level 1 (Topological).}  The quantities $\nu$,
    $\mufree$, the $S_N$ irrep, and the structure type depend only
    on the electron count $N$ and the total spin $S$.  They are
    independent of $Z$, the orbital angular momenta, and all
    details of the electron--electron interaction.  They are
    \emph{universal}.

  \item \textbf{Level 2 (Perturbative/Metric).}  The binding energy,
    ionization energy, electron configuration, and Hund's rules all
    depend on $Z$, the fine-structure constant $\alpha$, and the
    shape of the potential.  These quantities determine the
    \emph{chemistry} within a given topological class.
\end{itemize}

An important consequence: \emph{Hund's first rule is not a topological
effect.}  Since $\lambda_1 = 2$ for all spatial irreps with $S < N/2$,
the angular quantum number $\nu = N - 2$ is independent of the spin
value.  The preference for high-spin ground states (Hund's first rule)
enters through the perturbative Coulomb and exchange integrals at
Level~2, not through a reduction of $\nu$ at Level~1.

\subsection{The Dirac limit as metric singularity}
\label{sec:dirac_limit}

The Dirac equation for hydrogen-like atoms predicts that the 1s
binding energy diverges at $Z\alpha = 1$ ($Z \approx 137$ for a point
nucleus):
\begin{equation}
  E_{\mathrm{1s}}^{\mathrm{Dirac}}
    = mc^2 \left(\sqrt{1 - (Z\alpha)^2} - 1\right).
  \label{eq:dirac_1s}
\end{equation}
Does this ``Dirac dive'' have a topological signature on $S^{3N-1}$?

At $N = 137$: $\nu = 135$, $\mufree = 36{,}450$, and $\nu/(3N-2) =
0.3301$.  All formulas are perfectly smooth.  The excitation fraction
approaches $1/3$ from below with no singularity, no phase transition,
and no special structure at $N = 137$.

The Dirac instability is instead a \emph{metric} effect.  The
non-relativistic Fock projection yields the round metric on $S^3$.
Relativistic corrections warp this metric by a conformal factor
\begin{equation}
  \delta_{\mathrm{1s}} = \frac{(Z\alpha)^2}{1 - (Z\alpha)^2},
  \label{eq:metric_correction}
\end{equation}
which diverges at $Z\alpha \to 1$.  This ``pinch'' is a conical
singularity at the nuclear pole of $S^{3N-1}$---the geometric avatar
of the Dirac dive.  The topology ($S_N$ irrep, $\nu$, $\mufree$)
is oblivious; the metric (curvature, conformal factor) is singular.

Table~\ref{tab:metric} quantifies the transition from perturbative
to singular.

\begin{table}[h]
\centering
\caption{Relativistic metric correction $\delta_{1\mathrm{s}} =
(Z\alpha)^2/[1 - (Z\alpha)^2]$ for the innermost electron.  The
non-relativistic $\mufree$ formula becomes unreliable when
$\delta \gtrsim 1$.}
\label{tab:metric}
\begin{tabular}{rlrl}
\toprule
$Z$ & $(Z\alpha)^2$ & $\delta_{1\mathrm{s}}$ & Regime \\
\midrule
  1 & 0.00005 & 0.00005 & Non-relativistic \\
 29 & 0.045   & 0.047   & Mildly relativistic \\
 47 & 0.118   & 0.133   & Relativistic \\
 79 & 0.332   & 0.498   & Strongly relativistic \\
118 & 0.741   & 2.87    & Near-singular \\
137 & 0.999   & $\sim 10^3$ & Conical singularity \\
\bottomrule
\end{tabular}
\end{table}

\subsection{Three limits on the periodic table}

The periodic table terminates due to physics beyond the
non-relativistic $S_N$ representation theory:
\begin{enumerate}
  \item \textbf{Nuclear instability} ($Z \sim 126$): proton--proton
    Coulomb repulsion overwhelms the strong nuclear force.
  \item \textbf{Dirac instability, point nucleus} ($Z \sim 137$): the
    1s binding energy reaches $mc^2$.
  \item \textbf{Dirac instability, finite nucleus} ($Z \sim 170$):
    spontaneous $e^+e^-$ pair creation from the vacuum.
\end{enumerate}
The SO$(3N)$ topology predicts no limit: $\mufree = 2(N-2)^2$ is
defined for all $N$.  The physical end of the periodic table is set by
nuclear physics ($Z \lesssim 130$), not by electronic topology.

%% ========================================================================
%% VII. DISCUSSION
%% ========================================================================

\section{Discussion}
\label{sec:discussion}

\subsection{Connection to Paper 0's orbital counting}

Paper~0~\cite{loutey_paper0} derived the $2n^2$ orbital degeneracy
from the geometric packing of nodes on the discrete $S^3$ lattice.
The present result extends this: the $S_N$ representation theory on
$S^{3N-1}$ yields not just the orbital counting but the
\emph{chemical classification} of atoms.  Where Paper~0 derived the $2n^2$ counting from geometric packing,
the present analysis provides a group-theoretic framework for
understanding why filling those shells produces the periodic pattern
$\mathrm{C} \to \mathrm{D} \to \mathrm{E} \to \mathrm{B}$.

\subsection{Transition metals and Hund's rule}

The transition metals (Type~E) have partially filled $d$-shells with
ground-state spins $S \geq 1$.  The ``anomalous'' configurations of
Cr~($[\mathrm{Ar}]\,3d^5\,4s^1$, $S = 3$) and
Cu~($[\mathrm{Ar}]\,3d^{10}\,4s^1$, $S = 1/2$) are often cited as
evidence for exchange stabilization.  In the present framework, both
the ``expected'' and ``anomalous'' configurations of Cr have
$\nu = 22$ and $\mufree = 1012$---the topological Level~1 cannot
distinguish them.  The anomaly lives entirely at Level~2, in the
exchange pair count: the half-filled $3d^5\,4s^1$ configuration has
141~exchange pairs versus 136 for $3d^4\,4s^2$, a difference of~5
that tips the energetic balance.

This confirms that Hund's first rule is a \emph{perturbative} effect
(exchange stabilization in the Coulomb interaction) rather than a
\emph{topological} one (angular momentum reduction on $S^{3N-1}$).
More broadly, the actual chemistry---which electron configuration wins
energetically---is determined by exchange, correlation, and
orbital-dependent Coulomb effects that operate below the resolution of
the topological classification.  The Level~1 framework identifies
which configurations are \emph{topologically equivalent}; it cannot
determine which one is the ground state.

\subsection{Physical chemistry correspondences}

The structure type classification connects naturally to familiar
chemical trends:
\begin{itemize}
  \item \textbf{Ionization energy} correlates with the topological
    cost of surgery on $S^{3N-1}$.  Type~B (noble gas) has the highest
    surgery cost; Type~C (alkali) has the lowest.
  \item \textbf{Atomic radius} correlates with the location of the
    surgery: Type~C atoms have a weakly bound outer electron (large
    radius), while Type~B atoms are compact.
  \item \textbf{Electronegativity} increases across a period
    ($\mathrm{C} \to \mathrm{B}$), mirroring the transition from
    hierarchical to democratic irreps.
  \item \textbf{The periodic law} can be understood as the
    repetition of the $S_N$ irrep sequence whenever a new principal
    shell begins to fill.
\end{itemize}

\subsection{Fractal structure: core and valence at every scale}

A striking feature of the classification is its \emph{scale
invariance}.  The core--valence decomposition at each ionization step
produces the same Types~B/C/D structure at every level:
\begin{align*}
  \text{Na} &\to \text{Na}^+ + e^-
    & (\text{C} &\to \text{B} + \text{A}), \\
  \text{Na}^+ &\to \text{Na}^{2+} + e^-
    & (\text{B} &\to \text{C} + \text{A}).
\end{align*}
Each ionization peels off one column from the Young diagram,
converting one type into another in a predictable sequence.  This
self-similar structure is a direct consequence of the recursive
nature of $S_N \supset S_{N-1}$ representations.

%% ========================================================================
%% VIII. CONCLUSION
%% ========================================================================

\section{Conclusion}
\label{sec:conclusion}

We have shown that the group structure of the periodic table---the
classification of elements into noble gases, alkali metals, alkaline
earths, and open-shell atoms---follows from the $S_N$ representation
theory of fermionic wavefunctions on the hypersphere $S^{3N-1}$.  The
key results are:

\begin{enumerate}
  \item The ground-state angular quantum number $\nu = N - 2$ is
    universal for $S < N/2$, yielding the Pauli centrifugal cost
    $\mufree = 2(N-2)^2$.

  \item The shape of the $S_N$ spatial irrep classifies atoms into
    types (A/B/C/D/E) that recast the chemical group structure in
    representation-theoretic language.

  \item The periodic law corresponds to the repetition of the irrep
    sequence $\mathrm{C} \to \mathrm{D} \to \mathrm{E} \to
    \mathrm{B}$ within each period.

  \item Hund's first rule and configuration anomalies (Cr, Cu) are
    Level~2 perturbative effects, invisible to the Level~1 topology.

  \item The Dirac instability at $Z\alpha \to 1$ is a metric
    singularity on $S^{3N-1}$, not a topological one; the
    representation theory is smooth through all $N$.
\end{enumerate}

The approach separates the periodic table into a universal
\emph{topological skeleton} (Level~1: $\nu$, $\mufree$, structure
type) and a chemistry-specific \emph{metric flesh} (Level~2: $Z$,
$\alpha$, Coulomb integrals).  This separation provides a
group-theoretic framework for the periodic table's structure
(topology) while leaving the element-specific chemistry to
perturbative effects (metric).  The framework recasts, rather than
derives from first principles, the empirical periodic pattern---its
value lies in unifying the classification under a single
mathematical structure.

\begin{acknowledgments}
This work uses the hyperspherical formalism developed in the Geometric
Vacuum program (Papers~0--15).  The $S_N$ analysis was informed by the
classical references of Avery~\cite{Avery2000} and
Fock~\cite{Fock1935, Fock1954}.
\end{acknowledgments}

%% ========================================================================
%% REFERENCES
%% ========================================================================

\begin{thebibliography}{25}

\bibitem{Mendeleev1869}
D.~I.~Mendeleev,
\emph{Zh.\ Russ.\ Khim.\ Obshch.} \textbf{1}, 60 (1869).

\bibitem{Fock1935}
V.~Fock,
\emph{Z.\ Phys.} \textbf{98}, 145 (1935).

\bibitem{Fock1954}
V.~Fock,
\emph{Izv.\ Akad.\ Nauk SSSR, Ser.\ Fiz.} \textbf{18}, 161 (1954).

\bibitem{Avery2000}
J.~Avery,
\emph{Hyperspherical Harmonics and Generalized Sturmians}
(Kluwer, Dordrecht, 2000).

\bibitem{Weyl1946}
H.~Weyl,
\emph{The Classical Groups: Their Invariants and Representations}
(Princeton University Press, Princeton, 1946).

\bibitem{FultonHarris1991}
W.~Fulton and J.~Harris,
\emph{Representation Theory: A First Course}
(Springer, New York, 1991).

\bibitem{loutey_paper0}
J.~Loutey,
``Geometric packing and the universal constant,''
GeoVac Paper~0 (2025).

\bibitem{loutey_paper1}
J.~Loutey,
``Spectral graph theory for quantum chemistry,''
GeoVac Paper~1 (2025).

\bibitem{loutey_paper7}
J.~Loutey,
``The dimensionless vacuum: Schr\"odinger recovery from $S^3$,''
GeoVac Paper~7 (2026).

\bibitem{loutey_paper13}
J.~Loutey,
``The hyperspherical lattice: two-electron atoms as coupled channel
graphs,''
GeoVac Paper~13 (2026).

\bibitem{KnowlesHandy1984}
P.~J.~Knowles and N.~C.~Handy,
\emph{Chem.\ Phys.\ Lett.} \textbf{111}, 315 (1984).

\bibitem{Dirac1928}
P.~A.~M.~Dirac,
\emph{Proc.\ R.\ Soc.\ London A} \textbf{117}, 610 (1928).

\bibitem{Greiner1985}
W.~Greiner, B.~M\"uller, and J.~Rafelski,
\emph{Quantum Electrodynamics of Strong Fields}
(Springer, Berlin, 1985).

\bibitem{Pyykkoe2012}
P.~Pyykk\"o,
\emph{Phys.\ Chem.\ Chem.\ Phys.} \textbf{13}, 161 (2011).

\bibitem{Schwerdtfeger2015}
P.~Schwerdtfeger, L.~F.~Pashtedsky, A.~Pernpointner, and B.~Fricke,
``Relativistic effects in superheavy elements,''
in \emph{The Chemistry of Superheavy Elements}, 2nd ed.\
(Springer, Berlin, 2015).

\end{thebibliography}

\end{document}
